\documentclass{article}
\usepackage{graphicx}
\usepackage{algorithm}
\usepackage{algpseudocode}

\title{Sistema multi-agente con stati aggregati}
\author{Lorenzo Pecorari}
\date{November 2025}

\begin{document}

\maketitle

\section{Spazio degli stati}
Dati $n > 0$ nodi e $T > 0$ intervalli temporali, siano rispettivamente $I = \{1, 2, ..., n\}$ e $\Theta = \{0, 1, ..., T\}$ rispettivamente  l'insieme degli indici dei nodi del sistema e l'insieme degli intervalli considerati.  

L'ambiente che simula il sistema caratterizzato dai nodi che agiscono nei specifici timestep presenta uno spazio degli stati $S$ tale che:
\begin{itemize}
    \item $S = \{s_1^{env}, s_2^{env}, ..., s_t^{env}, ..., s_k^{env}\}$
    \item $s_t^{env} = (s_t^{env,1}, s_t^{env, 2}, ..., s_t^{env, i}, ..., s_t^{env, n}) = (s_t^{env, i})_{i\in I}$
    \item $s_t^{env, i} = (b_t^i, \tau_t^i, f_{t-1}^i, x_{t-1}^i, g_{t-1}^i, h_{t-1}^i)$
\end{itemize}

dove:
\begin{itemize}
    \item $b^i_t, \tau^i_t$ rappresentano lo \textit{stato fisico} del nodo
    \item $f^i_{t-1}, x^i_{t-1}, g^i_{t-1}, h^i_{t-1}$ rappresentano 
          l'\textit{azione precedente} del nodo, necessarie per calcolare 
          lo stato aggregato $m^i_t$ al timestep successivo
\end{itemize}

Ogni stato $s_t^{env, i}$ presenta il livello normalizzato di batteria del nodo $i$ al tempo $t$, la quantità normalizzata di timestep completati rispetto alla durata dell'episodio, il framerate adottato localmente, l'azione di offloading effettuata, l'indice del nodo con cui tenta di fare matching al tempo $t-1$ e inifine il framerate che richiede/accetta.

È possibile quindi considerare ogni stato $s_t^{env}$ come lo stato congiunto al tempo $t$.

\section{Spazio delle osservazioni}
Per questioni di complessità computazionale, tenere traccia di tutti gli stati di tutti gli agenti risulta essere oneroso per un singolo nodo. Al fine di ovviare a tale problematica, ad ogni istante temporale ogni agente osserva l'ambiente e considera solo una parte ridotta di tutte le informazioni contenute in ogni stato congiunto. Tali informazioni vengono considerate come $o_t^i$, l'osservazione dell'environment al tempo $t$ da parte del nodo $i$.\\

Sia $O^i$ lo spazio delle osservazioni dell'ambiente condotte dall'agente $i$ tale che:
\begin{itemize}
    \item $O^i = \{o_0^i, o_1^i, ..., o_t^i, ..., o_T^i\} = \{o_t^i\}_{t\in\Theta}$
    \item $o_t^i = (l_t^i, m_t^i)$
    \item $l_t^i = (b_t^i, \tau_t^i)$, informazione locale
    \item $m_t^i = (mb_t^i, m\tau_t^i, mf_t^i, mh_t^i, req_t^i, acp_t^i)$, informazione globale
\end{itemize}

Ogni osservazione $o_t^i$ permette al nodo $i$ di essere aggiornato sia su parametri locali come il suo livello normalizzato di batteria $b_t^i$ che l'avanzamento dell'episodio $\tau_t^i$, sia su informazioni globali come la media delle batterie dei nodi, l'avanzamento medio, i framerate medi locali e di offloading e quanti nodi rispettivamente richiedono offloading e quanti lo accettano.

\section{Spazio delle azioni}
A livello di environment è possibile definire lo spazio delle azioni come l'insieme della azioni congiunte dei vari agenti. Sia $A = \{a_0, a_1, ..., a_t, ..., a_k\}$ dove $a_t = (a_t^1, a_t^2, ..., a_t^i, ..., a_t^n) = (a_t^i)_{i \in I}$, in dove ogni azione $a_t^i$ definita come $a_t^i = (f_t^i, x_t^i, g_t^i, h_t^i)$ rappresenta l'azione $t-esima$ del nodo $i$. Le variabili presenti nell'azione sono:
\begin{itemize}
    \item $f_i^j \in [0, fps_{max}^j]$, framerate adottato per computazione locale\\
    \item $x_i^j = $\{ \begin{array}{lr}
        $0 , ricezione o invio non abilitati$ \\
        $1 , invio lavoro ad altro nodo$\\
        $2 , ricezione lavoro da altro nodo$
        \end{array}\\
    \item $g_i^j = $\{ \begin{array}{lr}
        $0$, $\space se \space x_i^j$ = 0 \\
        $k \neq j$, se $x_i^j$ = 1 indice del nodo con cui stabilire offloading\\
        \end{array}\\
    \item $h_i^j \in [0, fps_{max}^j]$, framerate adottato per computazione di offloading
\end{itemize}
A partire da $A$, è possibile definire lo spazio delle azioni riferito ad ogni agente $i$ del sistema: $A^i = \{a_0^i, a_1^i, ..., a_t^i, ..., a_k^i\}$, dove ogni azione $a_t^i$ è definita come in precedenza.

\subsection{Definizione variabili dell'osservazione globale}
Attraverso la definizione delle azioni che possono eseguire i vari agenti, è possibile esprimere con maggior rigore le variabili dell'informazione globale dell'osservazione degli agenti:
\begin{itemize}
    \item $mb_t^i = \frac{1}{n-1} \sum_{j \neq i}^n b_{t-1}^j$, media dei livelli di batteria dei nodi $j \neq i$
    \item $m\tau_t^i = \frac{1}{n-1} \sum_{j \neq i}^n \tau_{t-1}^j$, media dei timesteps dei nodi $j \neq i$
    \item $mf_t^i = \frac{1}{n-1} \sum_{j \neq i}^n f_{t-1}^j$, media dei framerate locali dei nodi $j \neq i$
    \item $mh_t^i = \frac{1}{n-1} \sum_{j \neq i}^n h_{t-1}^j$, media dei framerate di offloading dei nodi $j \neq i$
    \item $req_t^i = \sum_{j \neq i}^{n} \mathbf{1}(x_{t-1}^j = 1)$, numero di nodi $j \neq i$ che richiedono offloading
    \item $acp_t^i = \sum_{j \neq i}^{n} \mathbf{1}(x_{t-1}^j = 2)$, numero di nodi $j \neq i$ che accettano offloading
\end{itemize}

dove la funzione $\mathbf{1}(arg)$ restituisce $1$ se $arg$ risulta vero, $0$ altrimenti.


\subsection{Tabella Q}
La necessità di rendere l'agente informato sullo stato degli altri nodi senza cadere in una dimensione esponenziale della tabella, come definito in precedenza, porta a considerare valori aggregati e inerenti ai valori medi dei livelli di batteria, dei framerate locali, dei framerate di offloading e il numero di agenti che richiedono e che accettano offloading. Tale approccio permette di mantenere contenuta la crescita del numero di elementi presenti nella Q table. Al fine di semplificare la notazione, è possibile considerare lo spazio delle osservazioni $O^i$ come lo spazio degli stati $S^i$ per l'agente $i$. Ne deriva che ogni stato $s_t^i$ del nodo $i$ è esprimibile come $s_t^i = o_t^i = (l_t^i, m_t^i)$ senza alterare in alcun modo la definizione delle informazioni locali $l_t^i$ e globali $m_t^i$. La tabella in questione per ogni agente, quindi, sarà ottenibile dal prodotto cartesiano tra lo spazio degli stati e lo spazio delle azioni: $Q^i =S^i \times A^i $.\\

Come in precedenza, $l_t^i$ definisce il livello della batteria e la percentuale di timestep completati, con 11 valori possibili per entrambi: $b_t^i \in [0.0, 1.0], \tau_t^i \in [0.0, 1.0]$ .\\

Gloablmente, invece, in $m_t^i$ si hanno $mf_t^i$ e $mh_t^i$, i valori medi dei framerate adottati localmente da ogni nodi $j \neq i$. Questi hanno come limite superiore $F_t^i = max(fps_{max}^j)_{j\neq i}$ e $H_t^i = max(fps_{max}^j)_{j \neq i}$. Sulla base di tale ragionamento, si assume che i valori considerabili di $mf_t^i$ e $mh_t^i$ saranno rispettivamente $F_t^i+1$ e $H_t^i+1$. Per quanto riguarda il numero di nodi che hanno richiesto ed accettato offloading, $req_t^i$ e $acp_t^i$, questo sarà compreso tra 0 e n; il limite superiore del loro valore sarà quindi $n+1$ .Infine, per i valori medi delle batterie e dei timesteps, il limite superiore sarà uguale a quello delle relative variabili "locali" di $l_t^i$.\\

Si assume quindi che $|S^i| = (11 \cdot 11) \cdot (11 \cdot 11 \cdot F_t^i \cdot H_t^i \cdot (n+1) \cdot (n+1)) = 11^2 \cdot (11^2 \cdot (n+1)^2 \cdot (max(fps_{max}^j)_{j \neq i} +1)^2 = 11^4 \cdot (n+1)^2 \cdot (max(fps_{max}^j)_{j \neq i} +1)^2$ . \\

In maniera simile allo spazio degli stati, è possibile definire lo spazio delle azioni caratterizzato dagli elementi $a_t^i = (f_t^i, x_t^i, g_t^i, h_t^i)$, dove:
\begin{itemize}
    \item $f_t^i \in [0, fps_{max}^i]$, considerabili $fps_{max}^i+1$ valori differenti
    \item $x_t^i \in \{0, 1, 2\}$, considerabili 3 valori differenti
    \item $g_t^i \in [0, n]$, considerabili $n+1$ valori con $n$ il numero di nodi
    \item $h_t^i \in [0, fps_{max}^i]$, considerabili $fps_{max}^i+1$ valori differenti
\end{itemize}

Seguendo l'approccio applicato per lo spazio degli stati, è possibile calcolare la dimensione dello spazio delle azioni "locali" $A^i$: $|A^i| = (fps_{max}^i + 1) \cdot 3 \space \cdot (n+1) \space \cdot (fps_{max}^i + 1) = 3(fps_{max}^i + 1)^2(n+1)$ .\\

Dalla conoscenza di entrambe le dimensioni degli spazi che vengono mantenuti nella Q table da un aegnte $i$, si ricava che $|Q^i| = |S^i \times A^i| = |S^i| \cdot |A^i| = 11^4 \cdot (n+1)^2 \cdot (max(fps_{max}^j)_{j \neq i} +1)^2 \cdot 3(fps_{max}^i + 1)^2(n+1) = 3 \cdot 11^4 \cdot (max(fps_{max}^j)_{j \neq i} +1)^2 \cdot (fps_{max}^i + 1)^2 \cdot (n+1)^3$

\subsubsection{Scelta dell'azione}
Seguendo l'algoritmo $\epsilon-greedy$, si cerca di ottenere un trade-off tra exploration ed exploitation per l'agente, basandosi sul decay factor $\epsilon_{dec} \in[0, 1]$, sul valore iniziale $\epsilon_0$ e su quello finale $\epsilon_{fin}$. Per ogni timestep $t \in \Theta$, l'agente sceglie casualmente l'azione con probabilità $\epsilon$, mentre con probabilità $1 - \epsilon$ "effettua un ragionamento" in modo da ottimizzare la ricompensa.

\subsubsection{Osservazione dell'environment}
Ogni volta che l'agente osserva l'ambiente, l'informazione ritornata è parziale rispetto allo stato di cui tiene traccia l'ambiente ma permette comunque al nodo di avere la conoscenza dei suoi parametri locali e una serie di valori che di stimare la situazione degli altri nodi nel timestep $t \in \Theta$. Rispetto ad avere tutto lo stato congiunto, tale apaproccio permette di mantenere una conoscenza di tutti i nodi senza eccedere con la dimensione della tabella Q.

\subsubsection{Aggiornamento della tabella}
Ad ogni timestep $t \in \Theta$, l'agente sceglie un'azione $a_t^i$ sulla base dell'osservazione ottenuta dall'ambiente. Secondo la policy $\pi_{\theta}^i$, viene scelta un'azione $a_t^i = \pi_\theta^i(s_t)$, eseguita e collezionato il reward. La policy da seguire, in questo caso, prevede la scelta  dell'azione che massimizza il valore Q in corrispondenza di tale stato $s_t$: $a_t^i = argmax_{a \in A^i} Q(s_t^i, a)$.\\\\ A seguito di ciò, si procede con l'aggiornamento della tabella secondo:\\

$Q(s_t^i, a_t^i) \leftarrow (1-\alpha)Q(s_t^i, a_t^i) + \alpha(r_t^i + \gamma(max_{a_{t+1}^i} Q(s_{t+1}^i, a_{t+1}^i)))$ \\\\
con $\alpha \in (0, 1)$ come learning rate e $\gamma \in (0, 1)$ come discount factor. 

\section{Funzione di reward}
Per ogni azione effettuata in ogni timestep $t \in \Theta$, ogni agente riceve una ricompensa immediata derivante dalla combinazione dei tre diversi reward. Di fatto, ogni nodo  ottiene $R(t)^i = r_t^i = r_{frames, t}^i + r_{battery, t}^i + r_{cooperation, t}^i$, dove:

\begin{itemize}
    \item $r_{frames, t}^i$ è la ricompensa per il rispetto del constraint relativo al framerate massimo ammissibile;
    \item $r_{battery, t}^i$ premia l'agente sulla base del corretto utlizzo dell'energia a disposizione e di come viene spesa per effettuare le computazioni locali e di offloading;
    \item $r_{cooperation, t}^i$ restituisce il valore sulla bontà della cooperazione tra due nodi.
\end{itemize}
Seguono gli pseudocodici che definiscono la logica dei vari reward

\begin{algorithm}
    \State{// Verifica se la somma degli fps adottati localmente e quelli di offloading sono minori uguali a quelli massimi per il nodo}
    \caption{Algoritmo per ricompensa per gestione delle frame, $r_{frames, t}^i$}\label{alg:cap}
    \begin{algoerithmic}
    \If{($(f_t^i + h_t^i) \leq fps_{max}^i$)}
        \State{return 1}
        %\State{return $\frac{(f_t^i + h_t^i)}{fps_{max}^i}$}
    \Else
    \State{return 0}
    \EndIf
    \end{algoerithmic}
\end{algorithm}

\begin{algorithm}
    \caption{Algoritmo per ricompensa per gestione della batteria, $r_{battery, t}^i$}\label{alg:cap}
    \begin{algoerithmic}
    \State{// Verifica se $i$ non effettua offloading}
    \If{($x_t^i =0$ \land $((f_t^i \cdot E_{F, loc} \cdot \Delta t) < (E_{B,t}^i + E_{PV, t}^i))$)}
        \State{return 1}\\
        %\State{return $(f_t^i \cdot E_{F, loc} \cdot \Delta t)$/ $(E_{B,t}^i + E_{PV, t}^i)$}\\

    \State{// Se $i$ invia lavoro, l'energia a disposizione deve essere necessaria}
    \ElsIf{($x_t^i =1$ \land $(((f_t^i \cdot E_{F, loc} \cdot \Delta t) + (h_t^i \cdot E_{F, send} \cdot \Delta t))) < (E_{B,t}^i + E_{PV, t}^i))$)}
        \State{return 1}\\
        
    %\State{return $((f_t^i \cdot E_{F, loc} \cdot \Delta t) + (h_t^i \cdot E_{F, send} \cdot \Delta t)) / (E_{B,t}^i + E_{PV, t}^i)$}\\


    \State{// Se $i$ riceve lavoro, l'energia a disposizione deve essere necessaria}
    \ElsIf{($x_t^i =2$ \land $(((f_t^i \cdot E_{F, loc} \cdot \Delta t) + (h_t^i \cdot E_{F, recv} \cdot \Delta t))) < (E_{B,t}^i + E_{PV, t}^i))$)}
        \State{return 1}\\
        %\State{return $((f_t^i \cdot E_{F, loc} \cdot \Delta t) + (h_t^i \cdot E_{F, recv} \cdot \Delta t)) / (E_{B,t}^i + E_{PV, t}^i)$}\\

    \State{// Vincolo non soddisfatto, ricompensa nulla}
    \Else
    \State{return 0}
    \EndIf
    \end{algoerithmic}
\end{algorithm}

\begin{algorithm}
    \caption{Algoritmo per ricompensa per cooperazione tra nodi, $r_{cooperation, t}^i$}\label{alg:cap}
    \begin{algoerithmic}

    \State{// Se $i$ non effettua offloading e non designa nessuno per tale operazione}
    \If{($x_t^i =0$ \land $g_t^i = 0$)}
        \State{return 1}\\

    \State{// Se $i$ effettua offloading, invia lavoro, designa un altro nodo e tale nodo intende ricevere da i}
    \ElsIf{($x_t^i =1$ \land $g_t^i \neq 0 \land g_t^i \neq i \land x_t^{g_t^i} = 2 \land g_t^{g_t^i} = i$)}
    \State{return 1}\\

    \State{// Se $i$ effettua offloading, riceve lavoro, designa un altro nodo e tale nodo intende inviare a i}
    \ElsIf{(($x_t^i =2$ \land $g_t^i \neq 0 \land g_t^i \neq i \land x_t^{g_t^i}= 1 \land g_t^{g_t^i} = i$))}
    \State{return 1}\\

    \State{// Nessun vincolo soddisfatto, ricompensa nulla}
    \Else
    \State{return 0}
    \EndIf
    \end{algoerithmic}
\end{algorithm}


\end{itemize}

\end{document}
